\section{Los errores que m\'as puntos quitan}

\begin{enumerate}
\item \textbf{No dibujar la regi\'on.} Con el dibujo los l\'imites salen solos; sin \'el se pierden el inciso y los que dependen de \'el.
\item \textbf{L\'imites de afuera con variables.} Los de la integral exterior siempre son n\'umeros.
\item \textbf{Olvidar el intervalo} al escribir una marginal o condicional (\guillemotleft$0\le x\le1$; 0 en otro caso\guillemotright).
\item \textbf{Confundir} $P(X<a\mid Y=y_0)$ (se usa la densidad condicional) con $P(X<a\mid Y<b)$ (se usa $P(A\cap B)/P(B)$).
\item \textbf{No sustituir $k$} en los incisos siguientes.
\item \textbf{Afirmar independencia sin mostrarla.} Escribe $f_X$, $f_Y$, su producto y la comparaci\'on.
\item \textbf{Redondear a media operaci\'on} en covariancias. Usa fracciones.
\item \textbf{No resaltar el resultado.} Las instrucciones del examen lo piden: enci\'erralo.
\end{enumerate}

\textbf{Comprobaciones de 10 segundos:} toda tabla, marginal y condicional suma o integra 1; toda probabilidad est\'a
entre 0 y 1; toda variancia es positiva; $-1\le\rho\le1$; $\E(X)$ cae dentro del intervalo de $X$.

\section{Simulacro tipo examen (90 minutos)}

Mismo formato que tu segundo parcial: 5 problemas, 20 puntos cada uno. Resu\'elvelo completo antes de ver las soluciones.

\begin{enumerate}[label=\textbf{\arabic*.-},leftmargin=*]
\item El n\'umero de reclamaciones $X$ y de devoluciones $Y$ que recibe una tienda en un d\'ia tiene la funci\'on de probabilidad conjunta
\begin{center}
\begin{tabular}{c|ccc}
$x\backslash y$ & 0 & 1 & 2\\ \hline
1 & 0.10 & 0.05 & 0.05\\
2 & 0.15 & 0.20 & 0.05\\
3 & 0.05 & 0.15 & 0.20
\end{tabular}
\end{center}
a) Obtener las marginales. b) Calcular $P(X+Y\le3)$. c) Obtener la funci\'on condicional de $Y$ dado $X=2$ y $P(Y\ge1\mid X=2)$.
d) ?`Son independientes? e) Calcular el coeficiente de correlaci\'on.

\item Sea $\fxy(x,y)=k(x+2y)$ para $0\le x\le1$, $0\le y\le2$, y cero en otro caso.
a) Determinar $k$. b) Calcular $P(X+Y\le1)$. c) Obtener las marginales. d) Obtener $f_{Y\mid x}$ y calcular $P(Y<1\mid X=\tfrac12)$.
e) Determinar la curva de regresi\'on de $Y$ dado $x$.

\item Sea $\fxy(x,y)=kx$ para $0\le y\le x\le1$, y cero en otro caso.
a) Determinar $k$. b) Calcular $P(X+Y\le1)$. c) ?`Son independientes? d) Calcular la covariancia y el coeficiente de correlaci\'on.

\item Las variables $X$ y $Y$ tienen $\E(X)=4$, $\Var(X)=9$, $\E(Y)=-2$, $\Var(Y)=4$ y $\rho=-0.5$.
Para $W=2X-3Y+5$ calcular $\E(W)$ y $\Var(W)$. ?`Cu\'anto valdr\'ia $\Var(W)$ si fueran independientes?

\item Un dep\'osito tiene al inicio de la semana una cantidad $Y$ (en miles de litros) de la cual se vende $X$, con
$\fxy(x,y)=6x$ para $0\le x\le y\le1$. Obtener la cantidad promedio que queda en el dep\'osito al final de la semana.
\end{enumerate}

\newpage
\subsection*{Soluciones del simulacro}

\textbf{1.} a) $f_X=(0.20,\,0.40,\,0.40)$ en $x=1,2,3$; $f_Y=(0.30,\,0.40,\,0.30)$ en $y=0,1,2$.

b) Celdas $(1,0),(1,1),(1,2),(2,0),(2,1),(3,0)$: $0.10+0.05+0.05+0.15+0.20+0.05=\resp{0.60}$.

c) Rengl\'on $x=2$ entre $0.40$: $f_{Y\mid X=2}=(0.375,\,0.5,\,0.125)$; $P(Y\ge1\mid X=2)=\resp{0.625}$.

d) $\fxy(1,0)=0.10\ne f_X(1)f_Y(0)=0.06$. \resp{\text{No}}.

e) $\E(X)=2.2$, $\E(Y)=1$, $\E(X^2)=5.4$, $\E(Y^2)=1.6$, $\Var(X)=0.56$, $\Var(Y)=0.6$.
$\E(XY)=0.05+0.10+0.40+0.20+0.45+1.20=2.4$. $\Cov=2.4-2.2=0.2$.
$\rho=0.2/\sqrt{0.56\cdot0.6}=\resp{0.345}$.

\textbf{2.} a) $\int_0^2\!\int_0^1(x+2y)\,dx\,dy=\int_0^2(\tfrac12+2y)\,dy=5\Rightarrow\resp{k=\tfrac15}$.

b) $\tfrac15\int_0^1\!\int_0^{1-x}(x+2y)\,dy\,dx=\tfrac15\int_0^1\big[x(1-x)+(1-x)^2\big]dx=\tfrac15\int_0^1(1-x)\,dx=\resp{\tfrac1{10}}$.

c) $f_X(x)=\tfrac15\int_0^2(x+2y)\,dy=\resp{\tfrac{2x+4}{5}}$, $0\le x\le1$;\quad
$f_Y(y)=\tfrac15\int_0^1(x+2y)\,dx=\resp{\tfrac{1+4y}{10}}$, $0\le y\le2$.

d) $f_{Y\mid x}=\dfrac{(x+2y)/5}{(2x+4)/5}=\resp{\dfrac{x+2y}{2x+4}}$, $0\le y\le2$.
Con $x=\tfrac12$: $\int_0^1\dfrac{0.5+2y}{5}\,dy=\dfrac{0.5+1}{5}=\resp{0.3}$.

e) $\E[Y\mid x]=\int_0^2y\,\dfrac{x+2y}{2x+4}\,dy=\dfrac{2x+\frac{16}{3}}{2x+4}=\resp{\dfrac{3x+8}{3(x+2)}}$.

\textbf{3.} a) $\int_0^1\!\int_0^xkx\,dy\,dx=k\int_0^1x^2dx=\tfrac k3\Rightarrow\resp{k=3}$.

b) La regi\'on $x+y\le1$ dentro de $y\le x$ es el tri\'angulo con v\'ertices $(0,0)$, $(1,0)$, $(\tfrac12,\tfrac12)$. Conviene flecha horizontal:
para $0\le y\le\tfrac12$, entra en $x=y$ y sale en $x=1-y$:
\[\int_0^{1/2}\!\!\int_y^{1-y}3x\,dx\,dy=\tfrac32\int_0^{1/2}\big[(1-y)^2-y^2\big]dy=\tfrac32\int_0^{1/2}(1-2y)\,dy=\tfrac32\cdot\tfrac14=\resp{\tfrac38}\]

c) Soporte triangular: \resp{\text{no}}. Con marginales: $f_X=3x^2$, $f_Y=\tfrac32(1-y^2)$, y $f_Xf_Y\ne3x$.

d) $\E(X)=\int_0^13x^3dx=\tfrac34$, $\E(X^2)=\tfrac35$, $\Var(X)=\tfrac35-\tfrac9{16}=\tfrac3{80}$.
$\E(Y)=\int_0^1\tfrac32y(1-y^2)\,dy=\tfrac38$, $\E(Y^2)=\tfrac15$, $\Var(Y)=\tfrac15-\tfrac9{64}=\tfrac{19}{320}$.
$\E(XY)=\int_0^1\!\int_0^x3x^2y\,dy\,dx=\int_0^1\tfrac32x^4dx=\tfrac3{10}$.
$\Cov=\tfrac3{10}-\tfrac9{32}=\resp{\tfrac3{160}=0.01875}$;\quad
$\rho=\dfrac{3/160}{\sqrt{(3/80)(19/320)}}=\resp{0.397}$.

\textbf{4.} $\E(W)=2(4)-3(-2)+5=\resp{19}$. $\Cov=\rho\,\sigma_X\sigma_Y=-0.5(3)(2)=-3$.
$\Var(W)=4(9)+9(4)+2(2)(-3)(-3)=36+36+36=\resp{108}$. Independientes: $36+36=\resp{72}$.

\textbf{5.} $\E(Y-X)=\int_0^1\!\int_0^y6x(y-x)\,dx\,dy=\int_0^1\big[3x^2y-2x^3\big]_0^ydy=\int_0^1y^3dy=\resp{\tfrac14}$ (250 litros).

\section{Respuestas del autoexamen de tus apuntes}

\begin{center}
\begin{tabular}{c|cccccccccc}
Pregunta & 1 & 2 & 3 & 4 & 5 & 6 & 7 & 8 & 9 & 10\\ \hline
Inciso & B & A & B & B & B & D & C & E & D & A\\
Valor & $\tfrac1{500}$ & 0.03125 & $\tfrac x6-\tfrac1{12}$ & $\tfrac{2x}{y^2}$ & 0.8 & 1.88 & 0 & 8 & 13, 12.87 & $-0.408$
\end{tabular}
\end{center}
El procedimiento de cada una est\'a en los cuadros \guillemotleft Mecaniza\guillemotright{} de las secciones anteriores.

\section{Plan de repaso y hoja de bolsillo}

\textbf{Orden sugerido} (unas 5 horas en total): secciones 2 y 3 (1 h), secci\'on 4 con l\'apiz (1.5 h),
secciones 5 a 7 (1.5 h), simulacro con reloj (1.5 h incluyendo revisi\'on). Lo que falles en el simulacro, rep\'itelo al d\'ia siguiente.

\begin{teoria}{Todo el tema en una caja}
\begin{center}
\begin{tabular}{ll}
Hallar $k$ & $\iint\fxy=1$\\
Probabilidad & sumar celdas / integrar sobre la regi\'on dibujada\\
Marginal & sumar o integrar la \emph{otra} variable\\
Condicional & conjunta $\div$ marginal de lo conocido\\
Independencia & ?`$\fxy=f_Xf_Y$ en todo punto? (soporte no rectangular: no)\\
$\E[g(X,Y)]$ & $\sum\sum g\,\fxy$ \ o \ $\iint g\,\fxy$\\
Regresi\'on & $\E[Y\mid x]=\int y\,f_{Y\mid x}\,dy$\\
Covariancia & $\E(XY)-\E(X)\E(Y)$\\
Correlaci\'on & $\rho=\Cov/(\sigma_X\sigma_Y)$, \ $\rho^2$ = determinaci\'on\\
Combinaci\'on lineal & $\E=a\mu_X+b\mu_Y+c$; \ $\Var=a^2\sigma_X^2+b^2\sigma_Y^2+2ab\Cov$
\end{tabular}
\end{center}
\end{teoria}

\textbf{Calculadora TI-Nspire CX II.} Esta versi\'on no incluye todav\'ia las indicaciones de qu\'e pesta\~na del programa usar en cada
tipo de ejercicio. Se agregar\'an cuando el programa est\'e hecho, como acordamos.
